\section{Preadditive and additive categories}
\label{section-additive-categories}

\noindent
Here is the definition of a preadditive category.

\begin{definition}
\label{definition-preadditive}
A category $\mathcal{A}$ is called {\it preadditive} if each
morphism set $\Mor_\mathcal{A}(x, y)$ is endowed
with the structure of an abelian group such that the
compositions
$$
\Mor(x, y) \times \Mor(y, z)
\longrightarrow
\Mor(x, z)
$$
are bilinear. A functor $F : \mathcal{A} \to \mathcal{B}$ of
preadditive categories is called {\it additive} if and only
if $F : \Mor(x, y) \to \Mor(F(x), F(y))$
is a homomorphism of abelian groups for all
$x, y \in \Ob(\mathcal{A})$.
\end{definition}

\noindent
In particular for every $x, y$ there exists at least
one morphism $x \to y$, namely the zero map.

\begin{lemma}
\label{lemma-preadditive-zero}
Let $\mathcal{A}$ be a preadditive category.
Let $x$ be an object of $\mathcal{A}$.
The following are equivalent
\begin{enumerate}
\item $x$ is an initial object,
\item $x$ is a final object, and
\item $\text{id}_x = 0$ in $\Mor_\mathcal{A}(x, x)$.
\end{enumerate}
Furthermore, if such an object $0$ exists, then a morphism
$\alpha : y \to z$ factors through $0$ if and only if $\alpha = 0$.
\end{lemma}

\begin{proof}
First assume that $x$ is either (1) initial or (2) final.
In both cases, it follows that $\Mor(x,x)$ is a trivial abelian group
containing $\text{id}_x$, thus $\text{id}_x = 0$ in
$\Mor(x, x)$, which shows that each of (1) and (2) implies (3).

\medskip\noindent
Now assume that $\text{id}_x = 0$ in $\Mor(x,x)$. Let $y$
be an arbitrary object of $\mathcal{A}$ and let $f \in \Mor(x ,y)$.
Denote $C : \Mor(x,x) \times \Mor(x,y) \to \Mor(x,y)$ the composition map.
Then $f = C(0, f)$ and since $C$ is bilinear we have $C(0, f) = 0$.
Thus $f = 0$. Hence $x$ is initial in $\mathcal{A}$.
A similar argument for $f \in \Mor(y, x)$ can be used to show that
$x$ is also final. Thus (3) implies both (1) and (2).
\end{proof}

\begin{definition}
\label{definition-zero-object}
In a preadditive category $\mathcal{A}$, we call an object that is both
final and initial (as in Lemma \ref{lemma-preadditive-zero}) a
{\it zero object} and denote it $0$.
\end{definition}

\begin{lemma}
\label{lemma-preadditive-direct-sum}
Let $\mathcal{A}$ be a preadditive category.
Let $x, y \in \Ob(\mathcal{A})$.
If the product $x \times y$ exists, then so does
the coproduct $x \amalg y$.
If the coproduct $x \amalg y$ exists, then so does
the product $x \times y$. In this case
also $x \amalg y \cong x \times y$.
\end{lemma}

\begin{proof}
Suppose that $z = x \times y$ with projections
$p : z \to x$ and $q : z \to y$. Denote $i : x \to z$
the morphism corresponding to $(1, 0)$. Denote $j : y \to z$
the morphism corresponding to $(0, 1)$. Thus we have the
commutative diagram
$$
\xymatrix{
x \ar[rr]^1 \ar[rd]^i & & x \\
& z \ar[ru]^p \ar[rd]^q & \\
y \ar[rr]^1 \ar[ru]^j & & y
}
$$
where the diagonal compositions are zero. It follows that
$i \circ p + j \circ q : z \to z$ is the identity since
it is a morphism which upon composing with $p$ gives $p$
and upon composing with $q$ gives $q$.
Suppose given morphisms $a : x \to w$ and $b : y \to w$.
Then we can form the map $a \circ p + b \circ q : z \to w$.
In this way we get a bijection $\Mor(z, w)
= \Mor(x, w) \times \Mor(y, w)$ which
shows that $z = x \amalg y$.

\medskip\noindent
We leave it to the reader to construct the morphisms
$p, q$ given a coproduct $x \amalg y$ instead of a
product.
\end{proof}

\begin{definition}
\label{definition-direct-sum}
Given a pair of objects $x, y$ in a preadditive category $\mathcal{A}$,
the {\it direct sum} $x \oplus y$ of $x$ and $y$ is the direct
product $x \times y$ endowed with the morphisms
$i, j, p, q$ as in Lemma \ref{lemma-preadditive-direct-sum} above.
\end{definition}

\begin{remark}
\label{remark-direct-sum}
Note that the proof of Lemma \ref{lemma-preadditive-direct-sum}
shows that given $p$ and $q$ the morphisms $i$, $j$ are uniquely
determined by the rules $p \circ i = \text{id}_x$,
$q \circ j = \text{id}_y$, $p \circ j = 0$, $q \circ i = 0$.
Moreover, we automatically have
$i \circ p + j \circ q = \text{id}_{x \oplus y}$.
Similarly, given $i$, $j$ the morphisms $p$ and $q$ are uniquely determined.
Finally, given objects $x, y, z$ and morphisms
$i : x \to z$, $j : y \to z$, $p : z \to x$ and
$q : z \to y$ such that $p \circ i = \text{id}_x$,
$q \circ j = \text{id}_y$, $p \circ j = 0$, $q \circ i = 0$
and $i \circ p + j \circ q = \text{id}_z$, then $z$
is the direct sum of $x$ and $y$ with the four morphisms
equal to $i, j, p, q$.
\end{remark}

\begin{lemma}
\label{lemma-additive-additive}
Let $\mathcal{A}$, $\mathcal{B}$ be preadditive categories.
Let $F : \mathcal{A} \to \mathcal{B}$ be an additive functor.
Then $F$ transforms direct sums to direct sums and zero to zero.
\end{lemma}

\begin{proof}
Suppose $F$ is additive. A direct sum $z$
of $x$ and $y$ is characterized by having morphisms
$i : x \to z$, $j : y \to z$, $p : z \to x$ and
$q : z \to y$ such that $p \circ i = \text{id}_x$,
$q \circ j = \text{id}_y$, $p \circ j = 0$, $q \circ i = 0$
and $i \circ p + j \circ q = \text{id}_z$, according
to Remark \ref{remark-direct-sum}. Clearly $F(x), F(y), F(z)$
and the morphisms $F(i), F(j), F(p), F(q)$ satisfy exactly the
same relations (by additivity) and we see that $F(z)$ is
a direct sum of $F(x)$ and $F(y)$.
Hence, $F$ transforms direct sums to direct sums.

\medskip\noindent
To see that $F$ transforms zero to zero, use the
characterization (3) of the zero object in
Lemma \ref{lemma-preadditive-zero}.
\end{proof}

\begin{definition}
\label{definition-additive-category}
A category $\mathcal{A}$ is called {\it additive}
if it is preadditive and finite products exist, in other
words it has a zero object and direct sums.
\end{definition}

\noindent
Namely the empty product is a finite product and
if it exists, then it is a final object.

\begin{definition}
\label{definition-kernel}
Let $\mathcal{A}$ be a preadditive category.
Let $f : x \to y$ be a morphism.
\begin{enumerate}
\item A {\it kernel} of $f$ is a morphism
$i : z \to x$ such that (a) $f \circ i = 0$ and (b)
for any $i' : z' \to x$ such that $f \circ i' = 0$ there
exists a unique morphism $g : z' \to z$ such that
$i' = i \circ g$.
\item If the kernel of $f$ exists, then we denote
this $\Ker(f) \to x$.
\item A {\it cokernel} of $f$ is a morphism
$p : y \to z$ such that (a) $p \circ f = 0$ and (b)
for any $p' : y \to z'$ such that $p' \circ f = 0$ there
exists a unique morphism $g : z \to z'$ such that
$p' = g \circ p$.
\item If a cokernel of $f$ exists we denote this
$y \to \Coker(f)$.
\item If a kernel of $f$ exists, then a {\it coimage
of $f$} is a cokernel for the morphism $\Ker(f) \to x$.
\item If a kernel and coimage exist then we denote this
$x \to \Coim(f)$.
\item If a cokernel of $f$ exists, then the {\it image of
$f$} is a kernel of the morphism $y \to \Coker(f)$.
\item If a cokernel and image of $f$ exist then we denote
this $\Im(f) \to y$.
\end{enumerate}
\end{definition}

\noindent
In the above definition, we have spoken of ``the kernel'' and
``the cokernel'', tacitly using their uniqueness
up to unique isomorphism. This follows from the Yoneda lemma
(Categories, Section \ref{categories-section-opposite})
because the kernel of $f : x \to y$ represents the functor
sending an object $z$ to the set
$\Ker(\Mor_\mathcal{A}(z, x) \to \Mor_\mathcal{A}(z, y))$.
The case of cokernels is dual.

\medskip\noindent
We first relate the direct sum to kernels as follows.

\begin{lemma}
\label{lemma-additive-cat-biproduct-kernel}
Let $\mathcal{C}$ be a preadditive category.
Let $x \oplus y$ with morphisms $i, j, p, q$ as in
Lemma \ref{lemma-preadditive-direct-sum}
be a direct sum in $\mathcal{C}$. Then $i : x \to x \oplus y$
is a kernel of $q : x \oplus y \rightarrow y$. Dually, $p$ is
a cokernel for $j$.
\end{lemma}

\begin{proof}
Let $f : z' \to x \oplus y$ be a morphism such that $q \circ f = 0$.
We have to show that there exists a unique morphism $g : z' \to x$
such that $f = i \circ g$. Since $i \circ p + j \circ q$ is the identity on
$x \oplus y$ we see that
$$
f = (i \circ p + j \circ q) \circ f = i \circ p \circ f
$$
and hence $g = p \circ f$ works. Uniqueness holds because $p \circ i$
is the identity on $x$. The proof of the second statement is dual.
\end{proof}

\begin{lemma}
\label{lemma-kernel-mono}
Let $\mathcal{C}$ be a preadditive category.
Let $f : x \to y$ be a morphism in $\mathcal{C}$.
\begin{enumerate}
\item If a kernel of $f$ exists, then
this kernel is a monomorphism.
\item If a cokernel of $f$ exists, then
this cokernel is an epimorphism.
\item If a kernel and coimage of $f$ exist, then
the coimage is an epimorphism.
\item If a cokernel and image of $f$ exist, then
the image is a monomorphism.
\end{enumerate}
\end{lemma}

\begin{proof}
Part (1) follows easily from the uniqueness required in the
definition of a kernel. The proof of (2) is dual.
Part (3) follows from (2), since the coimage is a cokernel.
Similarly, (4) follows from (1).
\end{proof}

\begin{lemma}
\label{lemma-coim-im-map}
Let $f : x \to y$ be a morphism in a preadditive category
such that the kernel, cokernel, image and coimage all exist.
Then $f$ can be factored uniquely as
$x \to \Coim(f) \to \Im(f) \to y$.
\end{lemma}

\begin{proof}
There is a canonical morphism $\Coim(f) \to y$
because $\Ker(f) \to x \to y$ is zero.
The composition $\Coim(f) \to y \to \Coker(f)$
is zero, because it is the unique morphism which gives
rise to the morphism $x \to y \to \Coker(f)$ which
is zero
(the uniqueness follows from
Lemma \ref{lemma-kernel-mono} (3)).
Hence $\Coim(f) \to y$ factors uniquely through
$\Im(f) \to y$, which gives us the desired map.
\end{proof}

\begin{example}
\label{example-not-abelian}
Let $k$ be a field.
Consider the category
of filtered vector spaces over $k$.
(See Definition \ref{definition-filtered}.)
Consider the filtered vector spaces $(V, F)$ and $(W, F)$ with
$V = W = k$ and
$$
F^iV
=
\left\{
\begin{matrix}
V & \text{if} & i < 0 \\
0 & \text{if} & i \geq 0
\end{matrix}
\right.
\text{ and }
F^iW
=
\left\{
\begin{matrix}
W & \text{if} & i \leq 0 \\
0 & \text{if} & i > 0
\end{matrix}
\right.
$$
The map $f : V \to W$ corresponding to $\text{id}_k$ on the underlying
vector spaces has trivial kernel and cokernel but is not
an isomorphism. Note also that $\Coim(f) = V$ and $\Im(f) = W$.
This means that the category of filtered vector spaces over $k$
is not abelian.
\end{example}

\noindent
We finish this section by some lemmas on splitting images of
idempotent endormorphisms. Recall that an endomorphism $f$
is idempotent if $f \circ f = f$.

\begin{remark}
\label{remark-idempotent-symmetry}
Let $\mathcal{A}$ be a preadditive category,
$x$ an object of $\mathcal{A}$,
and $f : x \to x$ idempotent.
Writing $g = \text{id} _ x - f$,
one has $f \circ g = 0$, $g \circ f = 0$, and $g \circ g = g$.
\end{remark}

\begin{lemma}
\label{lemma-idempotent-kernel-cokernel}
Let $\mathcal{A}$ be a preadditive category,
$x$ be an object of $\mathcal{A}$,
and $f : x \to x$ be idempotent.
Write $g = \text{id} _ x - f$
and recall by Remark \ref{remark-idempotent-symmetry}
that $f \circ g = g \circ f = 0$.
\begin{enumerate}
\item If $f$ has kernel $i : \Ker (f) \to x$
and $p$ is the unique morphism such that $g = i \circ p$,
then $p \circ i = \text{id} _ {\Ker (f)}$
and $p$ is the cokernel of $f$.
\item If $f$ has cokernel $p : x \to \Coker (f)$
and $i$ is the unique morphism such that $g = i \circ p$,
then $p \circ i = \text{id} _ {\Coker (f)}$
and $i$ is the kernel of $f$.
\item If $g$ has kernel $j : \Ker (g) \to x$
and $q$ is the unique morphism such that $f = j \circ q$,
then $q \circ j = \text{id} _ {\Ker (g)}$
and $q$ is the cokernel of $g$.
\item If $g$ has cokernel $q : x \to \Coker (g)$
and $j$ is the unique morphism such that $f = j \circ q$,
then $q \circ j = \text{id} _ {\Coker (g)}$
and $j$ is the kernel of $g$.
\end{enumerate}
\end{lemma}

\begin{proof}
As for (1), compute first that $i = (f + g) \circ i = g \circ i$.
By the construction of $p$, this implies that
$i \circ p \circ i = g \circ i = i \circ \text{id} _ {\Ker (f)}$.
By Lemma \ref{lemma-kernel-mono} we conclude that
$p \circ i = \text{id} _ {\Ker (f)}$.

\medskip\noindent
It remains to show that $p : x \to \Ker (f)$
satisfies the universal property of the cokernel of $f$,
i.e., that given a morphism $a : x \to y$ satisfying $a \circ f = 0$
there exists a unique morphism $b : \Ker (f) \to y$
such that $a = b \circ p$.
Existence follows by first computing that $a = a \circ (f + g) = a \circ g$
and concluding by the construction of $p$
that $a = a \circ g = a \circ i \circ p$
(so that we can take $b = a \circ i$).
Uniqueness follows from that $b = b \circ p \circ i = a \circ i$
(so that $a$ determines $b$).

\medskip\noindent
By duality,
symmetry (i.e., that $g$ is itself idempotent;
cf. Remark \ref{remark-idempotent-symmetry}),
and duality and symmetry respectively, (2), (3), and (4) follow as did (1).
\end{proof}

\begin{lemma}
\label{lemma-idempotent-splitting}
Let $\mathcal{A}$ be a preadditive category,
$x$ be an object of $\mathcal{A}$,
and $f : x \to x$ be idempotent.
Write $g = \text{id} _ x - f$.
If each of $f$ and $g$ have at least one of a kernel or cokernel
and in turn morphisms $i , j , p , q$ are constructed
as in the statement of Lemma \ref{lemma-idempotent-kernel-cokernel},
then all four (co)kernels exist and
\begin{align*}
x
& \simeq \Ker (f) \oplus \Ker (g) \\
& \simeq \Coker (f) \oplus \Ker (g) \\
& \simeq \Ker (f) \oplus \Coker (g) \\
& \simeq \Coker (f) \oplus \Coker (g)
\end{align*}
with the structure of each direct product given by $i , j , p , q$.
\end{lemma}

\begin{proof}
By Lemma \ref{lemma-idempotent-kernel-cokernel}
$p \circ i = \text{id} _ {\Ker (f)} = \text{id} _ {\Coker (f)}$
and $q \circ j = \text{id} _ {\Ker (g)} = \text{id} _ {\Coker (g)}$.
By construction, $i \circ p + j \circ q = f + g = \text{id} _ x$.
By Remark \ref{remark-direct-sum} the claim follows.
\end{proof}

\begin{lemma}
\label{lemma-split-morphism-kernel-cokernel}
Let $\mathcal{A}$ be a preadditive category,
$x$ and $y$ be objects of $\mathcal{A}$,
and $j : y \to x$ and $q : x \to y$ be morphisms
satisfying $q \circ j = \text{id} _ y$.
Then
\begin{enumerate}
\item $\text{id} _ x - j \circ q$ has kernel $j$.
\item $\text{id} _ x - j \circ q$ has cokernel $q$.
\end{enumerate}
\end{lemma}

\begin{proof}
As for (1), compute first that $(\text{id} _ x - j \circ q) \circ j = 0$.
We must show that $j$ satisfies the universal property
of the kernel of $\text{id} _ {x} - j \circ q$,
i.e., that given a morphism $a : z \to x$
satisfying $(\text{id} _ x - j \circ q) \circ a = 0$
there exists a unique morphism $b : z \to y$
such that $a = j \circ b$.
Existence follows by computing that $a = j \circ q \circ a$
(so that we can take $b = q \circ a$).
Uniqueness follows from that $b = q \circ j \circ b = q \circ a$
(so that $a$ determines $b$).

\medskip\noindent
By duality, (2) follows as did (1).
\end{proof}

\begin{lemma}
\label{lemma-split-morphism-splitting}
Let $\mathcal{A}$ be a preadditive category,
$x$ and $y$ be objects of $\mathcal{A}$,
and $j : y \to x$ and $q : x \to y$ be morphisms
satisfying $q \circ j = \text{id} _ y$.
If $j \circ q$ has at least one of a kernel or cokernel,
then both (co)kernels exist and
\begin{align*}
x
& \simeq \Ker (j \circ q) \oplus y \\
& \simeq \Coker (j \circ q) \oplus y
\end{align*}
with the structure of either direct product including $j , q$,
and the canonical (co)kernel morphism,
and the remaining map uniquely determined.
\end{lemma}

\begin{proof}
Write $f = j \circ q$,
compute that $f \circ f = f$,
and apply Lemma \ref{lemma-idempotent-splitting}
and Lemma \ref{lemma-split-morphism-kernel-cokernel}.
\end{proof}




